While our focus is on single-sample hypothesis testing for a Bernoulli parameter,
the precision-goal framework extends naturally to other settings.
For single-group continuous data, Equation~\ref{eq:pitg_stop_iteration} generalises
by replacing $\theta(1-\theta)$ with the population variance $\sigma^2$; in practice
$\sigma^2$ is unknown and estimated by the sample variance $s^2$
(see Appendix~\ref{app:expected_stop} for the derivation).

For two-group comparisons --- whether binomial ($\theta_1 - \theta_2$) or continuous
($\mu_1 - \mu_2$) --- the standard error of the difference serves as the precision
metric, enabling sequential A/B testing under the same HDI+ROPE decision framework
\citep{kruschke2013}.
% \citet{kruschke2013}: Kruschke's BEST paper introduces the HDI+ROPE framework for
% two-group comparisons of proportions and means, making it the primary reference for
% extending the HDI+ROPE decision rule to the two-group setting discussed here.
For the binomial case with conjugate Beta posteriors, the posterior of the difference
can be evaluated by direct i.i.d.\ sampling from each Beta distribution separately,
without Markov Chain Monte Carlo (MCMC) \citep{gelman2013bayesian}.
% \citet{gelman2013bayesian}: BDA3 (Chapter 2/3) is the standard reference for
% deriving posterior quantities by direct Monte Carlo simulation. Cited here because
% it supports drawing i.i.d. samples from independent Beta posteriors to obtain the
% difference distribution exactly, without requiring a Markov chain.
For the continuous case, under a weakly informative prior the posterior of
$\mu_1 - \mu_2$ converges to the frequentist sampling distribution by the
Bernstein--von Mises theorem \citep{gelman2013bayesian}, so the standard error from
Welch's approximation provides a reliable surrogate for the posterior HDI width when
samples are of moderate size or larger.
% \citet{gelman2013bayesian}: BDA3 also covers the Bernstein--von Mises theorem and
% the conditions under which posteriors converge to the frequentist sampling distribution
% under weakly informative priors (Chapter 4), justifying the use of Welch's standard
% error as a surrogate for the posterior HDI width.

A practical advantage of the HDI+ROPE framework over frequentist approaches concerns
multiple comparisons. In NHST, conducting several tests simultaneously inflates the
Type~I error rate, requiring corrections (e.g.\ Bonferroni) that reduce the power of
each individual test. Frequentist confidence intervals inherit the same dependence on
the full set of intended comparisons, since their coverage properties are tied to
p-values. Consequently, both power and precision planning must account for the number
of planned tests.

The Bayesian posterior is not affected in this way. Because the posterior distribution
is determined solely by the observed data and the assumed data-generating process ---
not by the set of tests that were intended but not performed --- the power of any one
test is unaffected by how many other comparisons are planned, provided each test uses
an independent model \citep{gelman2012, kruschke2015doing}.
% \citet{gelman2012}: Gelman, Hill \& Yajima (2012) is the standard journal reference
% for the claim that Bayesian posteriors typically require no multiple-comparison
% correction when tests are based on independent models.
% \citet{kruschke2015doing}: DBDA2 (Chapter~11) provides the pedagogical treatment.
The precision goal $\omega_{\rm goal}$ and the stopping rule of ePitG therefore apply
unchanged in a multiple-testing context, without the need for correction.

This freedom from multiple-comparison penalties also reshapes how power analysis is
framed in the Bayesian setting.
It is worth distinguishing three types of power analysis that arise in this framework
\citep{kruschke2015doing}.

\textit{Prospective} power analysis is conducted before data collection: given a
hypothesised data-generating distribution (from theory, pilot data, or prior research),
one simulates experiments and asks how often the precision goal $\omega_{\rm goal}$ is
achieved within the budget $N_{\rm max}$.
This is precisely what the analytical formula for $N_{\rm goal}$
(Equation~\ref{eq:pitg_stop_iteration}) provides as a closed-form approximation for the
minimum required sample size within this model; note that $N_{\rm goal}$ is a theoretical
lower bound --- any sequential algorithm requires at least this many observations under
ideal conditions, and ePitG may require additional samples when $\theta_{\rm true}$ lies
near a ROPE boundary (see Section~\ref{sec:trends}).

\textit{Retrospective} power analysis uses the posterior from already-collected data as
the hypothetical data generator to ask whether the completed study was sufficiently
powered. In NHST this quantity is directly determined by the p-value and therefore adds
no new information beyond it \citep{hoenig2001}; in the Bayesian setting it can reveal
the probability of achieving goals other than null rejection.
% \citet{hoenig2001}: Hoenig \& Heisey (2001) demonstrate that observed power in NHST
% is a 1-to-1 function of the p-value and therefore adds no information beyond it ---
% the canonical reference for the redundancy of retrospective NHST power.
\textit{Replication} power asks: if the experiment were repeated exactly, what is the
probability of achieving the goal in the new study, informed by the results of the
first? This is straightforward to compute in a Bayesian framework by using the initial
posterior as the prior for the follow-up, but is difficult to operationalise in NHST,
which has no principled access to a posterior-derived data generator.

A recurring concern about PitG --- and by extension ePitG --- is that the precision
requirement feels too demanding in practice. 
Researchers frequently underestimate the noise in their measurements and are consequently surprised by the sample sizes that
$\omega_{\rm goal}$ implies (J.K.\ Kruschke, pers.\ comm., 2022). This is not a flaw
in the method but a form of epistemic honesty: imprecise posteriors produce inconclusive
results regardless of the algorithm applied; ePitG simply refuses to declare a verdict
until the posterior is genuinely informative. Setting $\omega_{\rm goal}$ thus forces
realistic planning rather than concealing underpowerment behind a false positive or a
silently inconclusive stop.
This discipline is a feature, not a limitation, for any analyst whose conclusions
must rest on genuinely informative data.

The method is most naturally suited to settings where data collection is relatively
cheap and the budget comfortably covers the required sample size: large-scale digital
A/B testing \citep{johari2022}, opinion polling, and automated quality-control pipelines
are typical examples.
% \citet{johari2022}: Johari et al.\ (2022) study continuous monitoring of A/B tests ---
% a canonical high-volume, low-cost data-collection context where precision-based
% stopping is operationally natural and the budget constraint is rarely binding.
For settings where the required sample size is prohibitive --- longitudinal studies,
clinical trials, or expensive policy evaluations --- the framework still serves a
valuable diagnostic role: $N_{\rm goal}$ makes explicit what precision is achievable
within the available budget, enabling an honest assessment of what conclusions the data
can and cannot support.

One option to increase signal with a limited budget is to incorporate informative priors.
We defer a full treatment of prior selection to future work, noting that informative
priors, where available and well-justified, offer a practical mitigation by contributing
effective pseudo-observations that narrow the starting HDI width and lower the required
sample size \citep{gelman2013bayesian}.
% \citet{gelman2013bayesian}: BDA3 (Chapter 2) establishes the prior-as-pseudo-data
% interpretation for conjugate models: a Beta(a,b) prior is equivalent to having
% previously observed a successes and b failures, directly supporting the claim that
% informative priors contribute effective pseudo-observations.

Three further considerations support the broader adoption of this approach.
First, formally accepting the null via the HDI+ROPE criterion parallels the objective
of frequentist equivalence testing through the Two One-Sided Tests (TOST) procedure
\citep{schuirmann1987, lakens2018}, providing a Bayesian counterpart to a methodology
already accepted in regulatory and applied settings \citep{fda2010bayesian}.
% \citet{fda2010bayesian}: The FDA's 2010 Bayesian guidance (\S4.7) explicitly permits
% stopping when the posterior credible interval is ``sufficiently narrow'', the same
% precision principle as PitG. ePitG adds the simultaneous conclusiveness requirement,
% directly addressing the inconclusive-outcome problem the FDA guidance leaves to
% case-by-case design agreement. Citing this document here supports the claim that
% precision-based Bayesian stopping is not merely academic but already embedded in
% regulatory practice.
Second, because $\omega_{\rm goal}$, the ROPE, and $N_{\rm max}$ must all be specified
before data collection begins, ePitG is naturally compatible with pre-registration
standards \citep{nosek2018}, directly reducing the researcher degrees of freedom that
drive selective reporting.
Third, the replication crisis has been partly attributed to underpowered studies whose
imprecise estimates fail to replicate \citep{osc2015, gelman2014}; ePitG addresses this
at the design stage by making precision a non-negotiable stopping criterion rather than
an aspiration --- one that is structurally enforced, since the stopping rule cannot fire
until the precision goal is met.
% \citet{osc2015}: The Open Science Collaboration (2015) documented that a substantial
% fraction of published psychological findings failed to replicate, motivating the search
% for design-stage remedies.
% \citet{gelman2014}: Gelman \& Carlin (2014) specifically link replication failure to
% low power and imprecision (``Type S'' sign errors and ``Type M'' magnitude errors),
% making it the canonical reference for the claim that imprecise estimates fail to replicate.

